% Simple poster (portrait)
% Author: Sofia Jijon (https://sjijon.github.io)
% Last Update: Sept 9, 2021
% Latest Version: https://github.com/sjijon/TeX-templates/tree/main/Tikzposter%20posters/Simple%20poster

\documentclass[a0paper,portrait,margin=0pt, colspace=24pt,subcolspace=0pt,blockverticalspace=36pt,innermargin=50pt]{tikzposter}

\usepackage[latin9]{inputenc}
\usepackage[square,numbers]{natbib} 	% Bibliography manager
\usepackage{amsmath,amssymb}
\usepackage{lipsum}  				    % Random Text
\usepackage[colalign]{aligncolsatbottom}  %To align columns at bottom (!! please run 2 times)

%..............................................................................................................................................................................................
% Display
\tikzposterlatexaffectionproofoff 			
\usetikzlibrary{shapes.geometric,arrows.meta,positioning}  %Tikz Libraries

% Fonts
\usepackage{helvet}					% Sans-Serif
\renewcommand{\familydefault}{\sfdefault}	%

% Colors
	\definecolor{MyOrange}{rgb}{0.8, 0.33, 0}
	\definecolor{MyBrown}{rgb}{0.28, 0.20, 0.20}
	\definecolor{MyGreen}{rgb}{0.33, 0.42, 0.18}

% Theme
\usetheme{Default}
\definecolorstyle{MyStyle2016}{
	\definecolor{ColorOne}{named}{MyBrown} 
	\definecolor{ColorTwo}{named}{MyOrange}
	\definecolor{ColorThree}{named}{MyGreen}
}{
    % Title Colors
    \colorlet{titlebgcolor}{ColorOne}
    \colorlet{titlefgcolor}{white}
    % Background Colors
    \colorlet{backgroundcolor}{ColorOne!15}
    \colorlet{framecolor}{ColorOne}
    % Block Colors
    \colorlet{blocktitlebgcolor}{white}
    \colorlet{blocktitlefgcolor}{ColorTwo}
    \colorlet{blockbodybgcolor}{white}
    \colorlet{blockbodyfgcolor}{black}
    % Innerblock Colors
    \colorlet{innerblocktitlebgcolor}{ColorOne!15}
    \colorlet{innerblocktitlefgcolor}{black}
    \colorlet{innerblockbodybgcolor}{ColorOne!15}
    \colorlet{innerblockbodyfgcolor}{black}
    % Note colors
    \colorlet{notebgcolor}{ColorTwo!20}
    \colorlet{notefgcolor}{ColorTwo}
    \colorlet{notefrcolor}{ColorTwo}
 }

% Color style
\usecolorstyle{MyStyle2016}
%..............................................................................................................................................................................................
\title{Modelling A Limb Congruency Learning Task}

\author{Antonio Amaddio}

\institute{FU-Berlin, Department of Psychology}

%..............................................................................................................................................................................................
\begin{document}
%
%
%	HEAD
%
%....................................................................................
%
%	Title
%
\maketitle[width=0.96\linewidth,titletoblockverticalspace=36pt,linewidth=0,roundedcorners=10]
%..............................................................................................................................................................................................
%
%	LEFT COLUMN
%
\begin{columns}
\column{0.33}
%....................................................................................
%
%	Block
%
\block[titleleft,roundedcorners=16]{1. Research Question}{
	\raggedright
	Is learning of motor actions (finger movement task) moderated by congruency in visual perception?
}
%....................................................................................
%
%	Block
%
\block[titleleft,roundedcorners=16]{2. Neural Correlates}{
	\raggedright
	Hypothesis 1: EBA activation is higher in congruent situation vs non congruent situation (in healthy patients)
}
%....................................................................................
%
%	Block
%
\block[titleleft,roundedcorners=16]{3. Task}{
	A finger movement task where (healthy) subjects perform better over time, naturally. Reward will be manipulated experimentally by visual correct/incorrect feedback.
	
	\tikzstyle{arrow} =  	      [->,outer sep=14pt, draw=ColorOne, text = ColorOne]
	\tikzstyle{enterexit} =      [rectangle, text width=1.5cm, text = ColorOne, draw=none]
	\tikzstyle{susceptible} =  [rectangle, text width=3cm, text = ColorThree!0, fill=ColorOne!95]
	\tikzstyle{infectious} =     [rectangle, text width=3cm, text = ColorTwo, fill=ColorTwo!25]
	\tikzstyle{recovered} =    [rectangle, text width=3cm, text = ColorThree!0, fill=ColorThree]

	\begin{center}
	\begin{tikzpicture}[node distance=6cm,text centered,minimum height=3cm, outer sep=4pt,inner sep=0pt,line width=6pt]
		\node (S) [susceptible] {\Large \textbf{S}};
		\node (I)  [right of=S,infectious] 	{\Large \textbf{I}};
		\node (R) [right of=I,recovered] 	{\Large \textbf{R}};
		
		\node (enterS) [above of=S,yshift=-1cm,enterexit] 	{ };
		\node (exitR) [below of=R,yshift=1cm,enterexit] 	{ };
		
		\draw [arrow] (enterS)  -- node[anchor=east] {$\pi$}  (S);
		
		\draw [arrow] (S)  -- 	node[anchor=south] {$\alpha$}  (I);
		\draw [arrow] (I)  --  node[anchor=south] {$\beta$}  (R);
		
		\draw[arrow] (R) -- node[anchor=east] { $\mu$} (exitR);				
	 \end{tikzpicture}
	 \end{center}
	 
 	That was a flowchart, and this is the corresponding ODE system:
	 
	 \innerblock[bodywidth=0.6\linewidth]{}{	
	\[
	\left\{
	\begin{aligned}
	    \dot{S} & =  \pi - \alpha S,\\ 
	    \dot{I}  & =   \alpha S - \beta I,\\
	    \dot{R} & = \beta I - \mu R. \\
	\end{aligned}
	\right.
	\]
	 }
}
%..............................................................................................................................................................................................
%
%	CENTER COLUMN
%
\column{0.34}
%....................................................................................
%
% 	Block
%
\block[titleleft,roundedcorners=16]{4. Hypothesis 1}{
	\raggedright
	This graphic was generated with Tikz:
	
	\begin{center}
	\begin{tikzpicture}[xscale=1.5,yscale=4,line width=4pt,text=ColorTwo]
		\draw [very thin, gray,step=0.5] (-pi,-1) grid (pi,1);
		\draw [color=ColorThree,domain=-pi:pi] plot (\x,{sin(\x r)});
		\draw (pi,0) node[right]{$\sin(x)$} ;
	\end{tikzpicture}
	 \end{center}

	\raggedright
	\lipsum[85]
	
	\begin{center}
	\begin{tabular}{l|c|c|r|}
		\textbf{Name} & \textbf{Age} & \textbf{Height [cm]} & \textbf{Color}\\
		\hline
		María & 25 & 152 & white\\
		Carlos & 12 & 142 & green\\
		Francisco & 18 & 172 & blue\\
		Mario & 21 & 165 & red
	\end{tabular}
	\end{center}
	
	\vspace{3cm}
	
	\lipsum[85]
 }
%....................................................................................
%
%	Block
%
\block[titleleft,roundedcorners=16]{5. Results H1}{
	\raggedright
	\lipsum[13]
}
%..............................................................................................................................................................................................
%
% 	RIGHT COLUMN
%
\column{0.33}
%%....................................................................................
%
%	Block
%
\block[titleleft,roundedcorners=16]{6. Hypothesis 2}{
	\raggedright
	
	\lipsum[85]
	\begin{center}
		\includegraphics[width=0.8\linewidth]{Figures/Image1}
	\end{center}
	
	}
%....................................................................................
%
%	Block
%
\block[titleleft,roundedcorners=16]{7. Results H2}{
	\raggedright
	\lipsum[12-13]
 }
\end{columns} 
%..............................................................................................................................................................................................
%
%	FOOT
%
%....................................................................................
%
%	References
%
\block[titleleft,roundedcorners=16]{}{
\small
\begin{minipage}{0.73\linewidth}
	\nocite{*}
	\bibliographystyle{unsrtnat}
	\bibliography{BibPoster}
 \end{minipage}
%....................................................................................
%
%	Logos
%
\begin{minipage}{0.2\linewidth}
\centering
	\includegraphics[height=5cm]{Figures/Logo_GitHub}
	\includegraphics[height=5cm]{github-link-qr-code}
\end{minipage}
}
%....................................................................................
%
%	My info
%
\note[width=14cm,targetoffsetx=3cm,targetoffsety=3cm,rotate=15]{
	\textbf{Contact information:}\\
	antonio.amaddio@fu-berlin.de
}
\end{document}